\documentclass[11pt]{article}
\usepackage[a4paper,margin=25mm]{geometry}
\usepackage{amsmath,amssymb,mathtools}
\usepackage{fontspec,unicode-math}
\setmainfont{Latin Modern Roman}
\setsansfont{Latin Modern Sans}
\setmathfont{Latin Modern Math}
\usepackage{xurl,hyperref}
\hypersetup{colorlinks=true,linkcolor=blue,urlcolor=blue,pdftitle={Problem statement}}
\newcommand{\sref}[2]{\hyperlink{source:#1:#2}{[S#2]}}
\newcommand{\cataloguescope}[1]{\paragraph{Mathematical question.}#1}
\begin{document}
\section{Moser's worm problem}\label{statement}

\subsection{Definitions and mathematical statement}

\paragraph{Shared notation.}$\mathbb N=\{1,2,\ldots\}$, and $\mathbb Z,\mathbb Q,\mathbb R,\mathbb C$ denote the integers, rationals, reals and complex numbers. Any other conventions are specified locally.


A worm is the image of a continuous rectifiable curve $\gamma:[0,1]\to\mathbb R^2$ of length one, where length is the supremum of sums $\sum_j|\gamma(t_j)-\gamma(t_{j-1})|$ over finite partitions. A planar region $K$ accommodates it if a translate and rotation of that image is contained in $K$. Area is Lebesgue area.
\paragraph{Statement recorded in the supplied source.}
Find a planar region of minimum area that accommodates every worm, and determine the minimum area. \sref{statement}{2}
\subsection{Short English statement}
Find the least-area planar region into which every continuous curve of length one can be moved and rotated.
\subsection{Sources}
\begin{description}
\item[\hypertarget{source:statement:1}{S1}] ULAM.AI and the UnsolvedMath Contributors, UnsolvedMath: A Curated Collection of Open Mathematics Problems, catalogue entry. CC BY 4.0. \url{https://huggingface.co/datasets/ulamai/UnsolvedMath}.
\item[\hypertarget{source:statement:2}{S2}] Zhipeng Deng, Universal Triangle Covering Curve and Polygonal Chain: Escaping Forest and Fitting Worm (2026), v2; arXiv:2608.01393. \url{https://arxiv.org/pdf/2608.01393}.
\end{description}
\label{end:statement}
\end{document}
